% !TEX root = ../main.tex
% =====================================================================
% 4. EXPERIMENTAL METHODOLOGY
%
% Only what a reader needs in order to judge the results. The full
% hyperparameter list, the commands and the environment are in
% Appendix C; the engineering details are kept short here.
% =====================================================================

\section{Experimental Methodology}
\label{sec:setup}

% ---------------------------------------------------------------------
\subsection{Datasets}
\label{sec:setup:data}

\begin{table}[h]
\caption{The two dataset families used in this study.}
\label{tab:datasets}
\begin{center}
\small
\begin{tabular}{lrrrl}
\toprule
\textbf{Dataset} & \textbf{\#series} & \textbf{Length} & \textbf{\#classes} &
\textbf{Per-step ground truth?} \\
\midrule
\webtraffic{}         & 1{,}000 (500 / 500) & 1{,}008 & 10 & yes \\
\ucr{} (85 published) & 40--16{,}637 & 24--2{,}709 & 2--60 & no \\
\ucr{} (128 total)    & 40--24{,}000 & 15--2{,}844 & 2--60 & no \\
\bottomrule
\end{tabular}
\end{center}
\end{table}
% Numbers checked against results/SEA_NET/data_summary.csv (train + test count per dataset) and
% each model's comparison_vs_millet.csv (the 85 rows with a non-empty millet_acc are the 85 MILLET
% published; the other 43 are the rest of the archive).

\webtraffic{} \citep{early2024millet} is synthetic, so its generator records which time steps it
modified to create each class -- the last column of \Cref{tab:datasets}, and the reason \webtraffic{}
is the only dataset on which NDCG@$n$ can be computed; on the \ucr{} archive \citep{dau2019ucr} only
AOPCR is available. \webtraffic{} is also small and fast to train on, so it was the screening set:
every new encoder and pooling head was evaluated there first and only the strongest combinations went
on to the full \ucr{} experiments, a choice whose cost \Cref{sec:res:ucr} quantifies.
\citet{early2024millet} published \millet{} results on 85 of the 128 \ucr{} datasets, so those 85
make a dataset-by-dataset comparison possible; the remaining 43 are run as well so that no conclusion
rests on one convenient subset.

% ---------------------------------------------------------------------
\subsection{Baselines and training configuration}
\label{sec:setup:baselines}

The baselines are FCN, ResNet \citep{wang2017fcn} and InceptionTime \citep{fawaz2020inceptiontime},
each paired with \millet{}'s conjunctive head. Their published accuracies are not quoted; all three
were re-trained through the same pipeline under the same training configuration as \seanet{}. This
was the most important methodological decision of the internship, because a published number combines
the effect of the architecture with that of the training configuration its authors used, and the two
cannot be separated afterwards. Re-training makes the comparison controlled: any difference between
\seanet{} and a re-trained baseline is attributable to the architecture. The size of the
configuration effect was measured rather than assumed -- a further run applies \millet{}'s own
published hyperparameters on this pipeline, and on \webtraffic{} the identical architecture scores
0.887 under the configuration of \Cref{tab:protocol} against 0.920 under theirs
(\Cref{sec:disc:reading}). Published \millet{} numbers therefore always occupy a separate row in
\Cref{sec:results}.

\label{sec:setup:protocol}
Every model, proposed and baseline alike, is trained with the identical configuration listed in
\Cref{tab:protocol}, so that differences in \Cref{sec:results} originate in the model. Two settings
adapt to the dataset rather than being fixed, because the \ucr{} archive spans 40 to 24{,}000
series: a 20\,\% validation split is held out only when a dataset has at least 100 training series,
below which the split would be too small to be informative and the training loss is monitored
instead; and the batch size scales with the training set size, since a fixed batch of 16 would
exceed half of the smallest \ucr{} datasets.

The number of seeds changed during the internship and the report distinguishes the two cases
throughout. The initial experiments were single-seed, because 72 variants over 129 datasets was
already close to the available compute budget (\Cref{app:constraints}); once the strongest models
were identified, two further seeds were added for the four models on which the main claims rest.
Variance estimates therefore exist exactly where the claims are made.

\begin{table}[h]
\caption{Training configuration, identical for every model in this study.}
\label{tab:protocol}
\begin{center}
\small
\begin{tabular}{ll}
\toprule
\textbf{Setting} & \textbf{Value} \\
\midrule
Optimiser              & Adam \citep{kingma2015adam}, lr $=1.25\times10^{-3}$, weight decay $=1.2\times10^{-4}$ \\
Epochs / early stopping& max 400 epochs; stop after 60 epochs with no improvement \\
Batch size             & $\min(\lfloor n_{\text{train}}/10 \rfloor, 16)$, floor of 2 \\
Validation split       & stratified 80/20 when $n_{\text{train}} \geq 100$; else train-loss early stopping \\
Label smoothing        & 0.13 \\
Focus penalty $\lambda_{\text{focus}}$ & 0.01 (\Cref{eq:loss}); 0 for the \millet{} baseline \\
Seeds                  & 3 for the four models in \Cref{tab:seeds}; 1 for the rest \\
Hardware               & Grid'5000 GPU nodes (Lille, Sophia) \\
\bottomrule
\end{tabular}
\end{center}
\end{table}

\paragraph{Metrics reported.}
Classification is measured by accuracy and test loss, explanation quality by AOPCR and NDCG@$n$
(\Cref{sec:bg:metrics}), and cost by the parameter count, the saved model size and the time for one
prediction. Cost is reported next to accuracy deliberately: without it, an improvement caused by a
better design cannot be distinguished from one caused by additional capacity. Over the \ucr{}
archive, results are summarised by the mean accuracy rank across datasets and by the win/tie/loss
record per dataset rather than by the mean accuracy alone, following the recommendation of
\citet{demsar2006statistical}, because a mean over heterogeneous datasets can be dominated by a
small number of them.

% ---------------------------------------------------------------------
\subsection{Experimental pipeline}
\label{sec:setup:pipeline}

\begin{contribution}
The experimental pipeline -- configuration system, model registry, resumable multi-dataset runs,
results store, and automatic figure and table generation -- is my own work. It is what made it
possible to train and compare 72 model variants over 129 datasets within the internship.
\end{contribution}

A model is always defined as \emph{encoder $+$ pooling head}, and both are selected by name in a
small configuration file rather than being hardcoded. Two registries, one for encoders and one for
heads, map a name to the corresponding class, and a builder connects the two. This is what allowed
72 variants to exist without 72 near-identical copies of the training code: every additional row in
the results tables is a new configuration file, not a new script, and each ablation of
\Cref{app:ablations} is a copy of an existing configuration with one line changed.

Two properties matter for the validity of the results. Each model writes into its own results
directory, one row per dataset per seed, so no two models can overwrite each other. And every figure
and \LaTeX{} table in this report is generated from those stored results by a single command
(\Cref{app:commands}), which is why no number reported here was transcribed by hand. The stack is
PyTorch \citep{paszke2019pytorch}, MLflow for tracking and Optuna \citep{akiba2019optuna} for the
search of \Cref{app:extra}. Development and small tests were local; every result reported here was
produced on Grid'5000 Linux GPU nodes. The remaining engineering details are in
Appendix~\ref{app:repro}.
